\subsection{Technical Appraisal and Evidence-Body Certainty}
\label{subsec:tqaf_results}

The deterministic TQAF appraisal covered all 206 included studies without
missing final scores. Overall evidence contribution was rated strong for 125
studies (60.7\%), adequate for 75 (36.4\%), and low for 6 (2.9\%). These
ratings describe the contribution that each study can make under the review's
reporting, validation, and claim-governance rules; they are not a binary
study-inclusion judgment.

Metric clarity was strong in 168 studies, adequate in 7, and low in 31.
Reporting completeness was strong in 196 and adequate in 10. In contrast, no
study reached the strongest benchmark-readiness category: 158 studies were
adequate and 48 were low. Comparison admissibility was strong in only 4
studies, adequate in 10, and low in 192. Reproducibility was strong in 3
studies, adequate in 199, and low in 4. This contrast indicates that complete
within-study reporting did not generally translate into portable benchmarks
or directly admissible cross-study comparisons.

Claim-level safeguards materially constrained the appraisal. The 31 studies
with at least one quarantined claim were capped at an overall score of 2. The
26 studies with a quarantined metric claim were capped at 1 for comparison
admissibility. Six studies combining low metric clarity with low validation
maturity were capped at an overall score of 1. In addition, 92 legacy metric
rows from 12 studies lacked a final comparability or admissibility label; these
were explicitly resolved as
\texttt{insufficient\_information\_due\_legacy\_extraction} rather than being
silently assigned a favorable score.

\begin{table*}[!t]
\caption{Study-level TQAF score distributions for the 206-study corpus.}
\label{tab:tqaf_distributions}
\centering
\begin{tabular}{lrrrr}
\toprule
\textbf{Dimension} & \textbf{Score 0} & \textbf{Score 1} & \textbf{Score 2} & \textbf{Score 3} \\
\midrule
Technical relevance & 0 & 15 & 68 & 123 \\
Metric clarity & 0 & 31 & 7 & 168 \\
Reporting completeness & 0 & 0 & 10 & 196 \\
Validation maturity & 0 & 32 & 168 & 6 \\
Reproducibility & 0 & 4 & 199 & 3 \\
Benchmark readiness & 0 & 48 & 158 & 0 \\
Comparison admissibility & 0 & 192 & 10 & 4 \\
Limitation transparency & 0 & 9 & 44 & 153 \\
Overall evidence contribution & 0 & 6 & 75 & 125 \\
\bottomrule
\end{tabular}
\end{table*}

The appraisal was also aggregated into 115 evidence bodies aligned with the
seven synthesis domains: 6 bodies for S1, 8 for S2, 47 for S3, 10 for S4, 3
for S5, 31 for S6, and 10 for S7. Certainty was high for 54 bodies, moderate
for 47, limited for 10, and unclear for 4. The four unclear bodies were broad
mixed or unclassified fallbacks retained to preserve coverage and
normalization provenance; they were explicitly designated non-substantive and
were not used to support survey conclusions.

\begin{table}[!t]
\caption{Evidence-body certainty across S1--S7.}
\label{tab:evidence_body_certainty}
\centering
\begin{tabular}{lr}
\toprule
\textbf{Certainty category} & \textbf{Evidence bodies} \\
\midrule
High & 54 \\
Moderate & 47 \\
Limited & 10 \\
Unclear/non-substantive fallback & 4 \\
\midrule
Total & 115 \\
\bottomrule
\end{tabular}
\end{table}

The appraisal therefore supports broad narrative conclusions about the
structure and technical direction of O-ISAC research, while requiring more
guarded language for benchmark-level performance comparisons. It also
identifies open artifacts, standardized baselines, and portable measurement
contracts---rather than reporting volume alone---as the principal constraints
on reproducible cross-study synthesis.
